\documentclass[11pt]{article}

\usepackage{amsmath, amssymb}
\usepackage{geometry}
\geometry{margin=1in}

\title{Information Distortion in Gaussian Predictive Inference}
\author{}
\date{}

\begin{document}

\maketitle

The Gaussian predictive model defines a local quadratic geometry of parameter space through the curvature of the total log-likelihood,
\begin{equation}
\mathbf H := -\mathbb{E}\!\left[\nabla^2_\theta \ell(\theta)\right],
\end{equation}
which we refer to as the Gaussian Fisher Information (GFI).

Independently, the total score
\begin{equation}
S(\theta) = \nabla_\theta \ell(\theta)
\end{equation}
fluctuates across stochastic realizations. Its covariance
\begin{equation}
\mathbf J := \mathrm{Cov}(S(\theta))
\end{equation}
defines the fluctuation geometry of the score process.

To compare these two geometries, we introduce the Information Distortion Matrix,
\begin{equation}
\mathbf C = \mathbf H^{-1/2} \mathbf J \mathbf H^{-1/2}.
\end{equation}

If the Gaussian predictive model fully decorrelates the measurement process,
the score behaves as an innovation process and $\mathbf C = \mathbf I$.
Deviations from identity quantify anisotropic distortion of information,
revealing residual temporal correlations or approximation errors.

Eigenvalues of $\mathbf C$ measure direction-dependent information mismatch,
while its determinant characterizes global information volume distortion.

\end{document}
